\documentclass{article}

\usepackage{sty/notes}
\usepackage{sty/math}

\title{CS 598 HPN Course Project Planning}
\author{Ian Chen}
\date{\today}

\usepackage[backend=biber,style=numeric,defernumbers=true]{biblatex}
\addbibresource{refs.bib}

\begin{document}

\maketitle

\section{Introduction}

My recent research has been about writing scalable graph algorithms ($\sim 5$ billion edges) for a single (shared-memory) machine.
However, these are limited by the memory and compute of a single machine.

In this project, I will write a communication substrate that brings the shared-memory programming model to distributed graph processing.
All operations are synchronous, with concurrent non-local memory reads handled via RDMA communication.

\subsection{Background}

In the shared memory case, the common programming model is to write graph algorithms as vertex or edge programs.
For example, consider the case of counting the number of triangles in a graph, we might write the formula

$$
    \sum_{uv \in E} \abs{N(u) \cap N(v)}
$$

This naturally expresses itself in code as a program to execute over the edges of a graph; this model naturally describes many other graph algorithms, including connectivity, shortest paths, or $k$-core.
However, synchronous algorithms for distributed graphs tend to be written the bulk-synchronous-parallel (BSP) model, which has a global barrier separating a communication phase before proceeding with local computation (as lots of small messages could make network communication the bottleneck)~\cite{dathathri2018gluon,steil2023embracing}.
The BSP model requires global synchronization (which amplifies load imbalance) and requires significantly more memory usage (to buffer messages into epochs)~\cite{mohammadiporshokooh2026overcoming}.
Moreover, some algorithms require many rounds, and these rounds can be heavily imbalanced in the work required, reducing performance.
Algorithms for asynchronous communication models have also been studied, but are work inefficient as they typically have to run until a convergence condition.
Both these alternatives requires a different algorithmic paradigm than is natural for a shared-memory machine, costing developer overhead.

Recently, kernel-bypass networking reduces the overhead of small messages.
As such, latency hiding has been proposed as an alternative~\cite{nelson2014grappa,mohammadiporshokooh2026overcoming}.
Edge operations access remote data directly and synchronously; these operations are run concurrently to utilize CPU.

\subsection{Evaluation}

In this project, I plan to write a communication substrate allowing the same program to be performant on one machine or on a cluster.
To minimize the latency, I plan to use LCI~\cite{yan2025lci} as a low-overhead RDMA communication layer.
On top, concurrent lightweight tasks hide the remaining latency.

The first experiment fixes the algorithm and graph and varies three conditions: the number of lightweight tasks per core, the network stack (LCI over RDMA, MPI, and TCP), and the number of machines.
We hypothesize that over RDMA, a few dozen tasks per core suffice to hide remote-access latency, keeping the same program within a small factor of Gluon as machines are added, while over TCP it falls far behind.
The second experiment fixes these conditions and varies the algorithm, comparing workloads where round-based algorithms converge slowly (connectivity on high-diameter graphs) against ones where rounds work well (PageRank).

\section{Related work}

Gluon~\cite{dathathri2018gluon} is a communication substrate that makes use of invariants of a static partitioned graph for optimizing communication.
However, the programming API is for the BSP model, and as such, suffers from the issues discussed before (namely, latency and developer overhead).

Recently, the NWGraph library has been combined with the HPX runtime system to reduce latency of distributed graph algorithms~\cite{mohammadiporshokooh2026overcoming}.
It supports lightweight task-based parallelism at a machine to hide the remaining latency of HPX.
Code must still be restructured to be aware of the placement of the data and place synchronization barriers.

\section{Computing resources}

CloudLab.

\printbibliography

\end{document}
